\documentclass[11pt]{scrartcl}
\usepackage[von,sexy]{evan}
\usepackage{booktabs}
\usepackage{wrapfig}

\begin{document}
\title{Curves of degree greater than 2}
\subtitle{MOP 2025}
\date{June 2025}
\maketitle

\section{Cayley-Bacharach}
\subsection{Curves of degree $d$}
We work in $\mathbb{CP}^2$ always in this lecture.
We'll often make use of the following.
\begin{lemma}
  [Bezout theorem]
  Two curves of degree $a$ and $b$ sharing no irreducible component
  intersect at up to $ab$ points (in fact exactly $ab$ ``with multiplicity'').
\end{lemma}
Then a degree-$d$ curve is given by a homogeneous polynomial of degree $d$ in
$\mathbb C[x,y,z]$, up to scaling.
This means a generic choice of $\binom{d+2}{2}-1$ points
should determine unique degree-$d$ curve: $2$ points determine a line,
$5$ points determine a conic, $9$ points determine a cubic, etc.
And for linear algebra reasons!

However, these are all only true in ``sufficiently general position''.
For example, if one takes four distinct collinear points,
any conic passing through three of them passes through the fourth.
Let's fix some (non-standard) terminology for this:
\begin{definition}
  Fix a degree $d \ge 0$.
  A set of points in $\mathbb{CP}^2$ is \emph{redundant} for degree-$d$ curves
  if vanishing on some proper subset implies vanishing on all of them.
\end{definition}
For example, three points are redundant for lines if and only if they're collinear.
(But four noncollinear points are redundant for lines too, due to vacuous truth ---
passing through $3$ points ``implies'' any conclusion.)

\subsection{Chasles theorem for cubics}
According to \cite{EGH}, this should be called Chasles theorem (due to Michel Chasles).
\begin{theorem}
  [Chasles]
  Suppose two plane cubics meet at exactly $9$ points. Then:
  \begin{enumerate}[(a)]
    \itemsep1pt
    \ii This set of $9$ points is redundant for cubics.
    \ii No subset of size $8$ is redundant.
    \ii Hence, any cubic passing through any $8$ points must pass through the ninth.
  \end{enumerate}
\end{theorem}
\noindent Part (c) is the one quoted by contestants, and follows from (a) and (b).
Example usage:
\begin{example}
  Deduce Pascal theorem from Chasles theorem.
\end{example}

\noindent Meanwhile, part (a) is obvious:
cubics are supposed to be determined by $9$ points,
but we're \emph{given} two distinct cubics pass through these.
The bulk of the work of the theorem is to prove (b),
and it turns out you can do this with just Bezout, no AG machinery:
\begin{problem}
  [{\cite[Proposition 1]{EGH}}]
  \label{prob:elem}
  Let $d \ge 1$ and consider $n$ points.
  Restrict attention\footnote{I
    want to emphasize $n \le 2d+2$ is a huge assumption.
    Removing this hypothesis is the main difficulty
    to general Cayley-Bacharach statements.}
  to the case $n \le 2d+2$.
  Show that these points are redundant for degree-$d$ curves if and only if
  either $d+2$ are collinear, or $n = 2d+2$ and all points lie on a conic.
\end{problem}

\begin{remark}
  [Elliptic curve addition]
  For experts: another application of Chasles
  is to show that addition on an elliptic curve is associative.
  In fact one can even use this on cubic curves for competition problems
  --- by defining a group law on the cubics to gain additional collinearities.
\end{remark}

\subsection{Cayley-Bacharach for all degrees}
Higher-degree versions of Cayley-Bacharach come in different forms.
Here is one formulation that I've adapted for contest use from \cite[Theorem CB5]{EGH};
however, it doesn't handle multiplicity.
(There are versions of Cayley-Bacharach that do handle multiplicity,
but they require language about schemes to state correctly.)
\begin{theorem}
  [Cayley-Bacharach]
  Suppose two curves of degree $a$ and $b$ meet at exactly $ab$ points,
  partitioned into two sets $S \sqcup T$.
  Let $d,e \ge 0$ with $d+e = a+b-3$.
  Then the following statements are equivalent:
  \begin{enumerate}[(i)]
    \itemsep1pt
    \ii Any degree-$d$ curve passing through all of $S$ passes through all of $T$.
    \ii The set $T$ is not redundant for curves of degree $e$.
  \end{enumerate}
\end{theorem}
\noindent Special cases:
\begin{itemize}
  \itemsep1pt
  \ii Chasles theorem is the special case where $a=b=3$, $d=3$, $e=0$, $|S|=8$, $|T|=1$.

  \ii Another common case called ``\alert{quartic Cayley-Bacharach}''\footnote{The
    projective geometry chapter of \cite{MODA}
    claims that the following condition on $S$ also works
    (instead of $T$ non-collinear):
    no five points of $S$ are collinear,
    no nine are co-conic, no twelve are co-cubic.
    Proving this seems to need a stronger non-elementary version
    of \Cref{prob:elem} without the condition that $n \le 2d+2$.
    In any case, contestants are often sloppy with hypotheses.
    I was careful to show Chasles theorem with 3 parts,
    to emphasize that higher-degree versions are more subtle.}
  by contestants is when $a=b=4$,
  $d=4$, $e=1$, $|S|=13$, and $|T|=3$ (note this requires $T$ non-collinear).
  Then any quartic passing through the first $13$ points passes through the last $3$.
\end{itemize}

\section{Making Cayley-Bacharach practical for competing}
\subsection{Circle points (most important and used all the time)}
In $\mathbb{CP}^2$, we consider two special points $\alpha$ and $\bar\alpha$
that every circle passes through, known as the \alert{circle points}.
In Cartesian coordinates they are defined via:
\begin{align*}
  \alpha &\coloneq (1:i:0) \\
  \bar{\alpha} &\coloneq (1:-i:0).
\end{align*}
\vspace{-0.7cm} % ???
\begin{definition}
  A cubic curve is \alert{circular} if it passes through $\alpha$ and $\bar\alpha$.
\end{definition}
\begin{exercise}
  Show that $\alpha$ and $\bar\alpha$ are isogonal conjugates
  with respect to any triangle.
  % The line at infinity swaps with the circumcircle, so it's clear
\end{exercise}
\begin{example}
  [Radical axis theorem]
  Show that the three common chords of any three pairwise intersecting circles concur.
\end{example}

\subsection{Isopivotal cubics (used sometimes)}
In what follows $ABC$ is a fixed triangle and we use barycentric coordinates.
\begin{definition}
  Let $ABC$ be a triangle and $P = (u:v:w)$ a point.
  The \alert{isopivotal cubic} with pivot $P$ is the cubic curve
  \[ 0 = \det \begin{bmatrix} u & v & w \\
    x & y & z \\ a^2 yz & b^2 zx & c^2 xy \end{bmatrix}.  \]
  It contains all $X$ such that $X$ and its isogonal conjugate
  $X^\ast$ are collinear with $P$.
\end{definition}
\begin{proposition}
  [Points on isopivotal cubics]
  A isopivotal cubic is self-isogonal and passes through all of the following points:
  the vertices $A$, $B$, $C$ of the triangle;
  the incenter and all three excenters;
  the pivot point $P$;
  the points $AP \cap BC$, $BP \cap CA$, $CP \cap AB$.
\end{proposition}
\begin{exercise}
  For which $P$ is the isopivotal cubic with pivot $P$ a circular cubic?
\end{exercise}

\noindent Many choices of $P$ have names.
Most prominently,
the \emph{Thomson cubic} (also called the 23-point cubic) has pivot $G$ (centroid),
and passes through $I$, $G$, $O$, $H$, $K$ and all altitude midpoints.
The \emph{Neuberg cubic} (also called the 21-point or 37-point curve)
has pivot $X_{30}$, the Euler infinity point (hence is circular),
the reflections of $A$, $B$, $C$ across the opposite sides,
and the six vertices of equilateral triangles constructed on the sides of $ABC$.
The \emph{McCay cubic} has pivot $O$ (circumcenter);
the \emph{orthocubic} has pivot $H$ (orthocenter);
the \emph{Darboux cubic} has pivot at the de Longchamps point.

\subsection{Isoptic cubics (used sometimes)}
Now we assume $ABCD$ is a quadrilateral.
It could be self-intersecting, but we assume no three of vertices are collinear,
and $ABCD$ is not a parallelogram.
\begin{theorem}
  [{\cite[Theorem 1.5]{Hu}}]
  The following are equivalent for a point $P$:
  \begin{enumerate}[(i)]
    \itemsep1pt
    \ii $\dang APB = \dang DPC$.
    \ii $\angle APC$ and $\angle BPD$ have the same bisectors.
    \ii The feet from $P$ to the sides are concyclic.
    \ii $P$ has an isogonal conjugate with respect to $ABCD$.
    \ii $P$ is the focus of some ellipse tangent to the four sides.
  \end{enumerate}
\end{theorem}
\begin{theorem}
  [{\cite[Definition 2.3; Theorem 7.1]{Hu}}]
  The set of such $P$ is a nondegenerate \alert{circular} cubic curve
  passing through the vertices of $ABCD$,
  called the \alert{isoptic cubic}\footnote{\cite{Hu}
    uses the name ``isogonal cubic'' instead.}
  of $ABCD$.
  The isoptic cubic passes through the Miquel point $M$ of $ABCD$;
  moreover, $M\alpha$ and $M\bar\alpha$ are tangent to the isoptic cubic.
\end{theorem}

\begin{thebibliography}{9}
\small
\bibitem[EGH95]{EGH}
  David Eisenbud, Mark Green, Joe Harris.
  \textit{Cayley-Bacharach theorems and conjectures.}

\bibitem[Gr12]{MODA}
  Adam Groucher. \textit{Mathematical Olympiad Dark Arts}.

\bibitem[Hu19]{Hu}
  Daniel Hu. \textit{Constructions in the locus of isogonal conjugates of a quadrilateral}.
  arXiv:1912.08296.
\end{thebibliography}

\section{Problems}
\subsection{Problems for class}
These problems collectively cover the main techniques for the lecture
without being too complicated.
If this material is mostly new to you, these are good ones to start with.

% vanilla Chasles application
\begin{problem}
  [Miquel point]
  Triangle $DEF$ is inscribed in triangle $ABC$.
  Show that the circumcircles of $AEF$, $BFD$, $CDE$ meet at a point.
\end{problem}

% isoptic cubic
\begin{problem}
  [Romania TST 2017]
  Let $ABCD$ be a convex quadrilateral and let $P$ and $Q$ be variable points
  inside this quadrilateral so that $\angle APB=\angle CPD=\angle AQB=\angle CQD$.
  Prove that the lines $PQ$ obtained in this way all pass through a fixed point,
  or they are all parallel.
\end{problem}

% degree-four Cayley Bacharach
\begin{problem}
  % https://artofproblemsolving.com/community/c6h515448p2897891
  % Groucher
  Let $ABCD$ be a convex cyclic quadrilateral whose diagonals meet at $E$.
  Prove that there exists a conic passing through the incenters of triangles
  $ABC$, $BCD$, $CDA$, $DAB$, $ABE$, $BCE$, $CDE$, $DAE$.
\end{problem}

% isopivotal cubic using the infinity point along line $AE$
\begin{problem}[\href{https://aops.com/community/p27349354}{USAMO 2023/6}]
Let $ABC$ be a triangle with incenter $I$
and excenters $I_a$, $I_b$, $I_c$ opposite $A$, $B$, and $C$, respectively.
Given an arbitrary point $D$ on the circumcircle of $\triangle ABC$
that does not lie on any of the lines $II_a$, $I_bI_c$, or $BC$,
suppose the circumcircles of $\triangle DII_a$ and $\triangle DI_bI_c$
intersect at two distinct points $D$ and $F$.
If $E$ is the intersection of lines $DF$ and $BC$,
prove that $\angle BAD = \angle EAC$.
\end{problem}

\begin{wrapfigure}[12]{r}{6cm}
\centering
\begin{asy}
/*
Converted from GeoGebra by User:Azjps using Evan's magic cleaner
https://github.com/vEnhance/dotfiles/blob/main/py-scripts/export-ggb-clean-asy.py
*/
pair B = (10.12122,2.49871);
pair A = (2.99463,1.73330);
pair C = (11.69783,-4.81262);
pair D = (2.18048,-4.59861);
pair E = (4.32162,1.87582);
pair F = (8.72099,2.34833);
pair G = (10.46541,0.90259);
pair H = (11.18169,-2.41907);
pair I = (9.88187,-4.77178);
pair J = (3.81855,-4.63544);
pair K = (2.58210,-1.47503);
pair L = (2.69159,-0.62352);
pair W = (4.15065,-0.33709);
pair X = (9.00350,0.61559);
pair Y = (9.46760,-2.23090);
pair Z = (4.05027,-1.63620);

import graph;
size(5cm);
pen zzttqq = rgb(0.6,0.2,0.);
pen qqwuqq = rgb(0.,0.39215,0.);
pen cqcqcq = rgb(0.75294,0.75294,0.75294);
draw(A--B--C--D--cycle, linewidth(0.6) + zzttqq);

draw(circle((7.,-2.), 5.47544), linewidth(0.6) + blue);
draw(A--B, linewidth(0.6) + zzttqq);
draw(B--C, linewidth(0.6) + zzttqq);
draw(C--D, linewidth(0.6) + zzttqq);
draw(D--A, linewidth(0.6) + zzttqq);
draw(circle((6.92,-1.6), 4.33969), linewidth(0.6) + blue);
draw(F--I, linewidth(0.6));
draw(E--J, linewidth(0.6));
draw(L--G, linewidth(0.6));
draw(H--K, linewidth(0.6));
draw(circle((6.83965,-1.19828), 2.82354), linewidth(0.6) + dashed + qqwuqq);

dotfactor *= 2;
dot(A); dot(B); dot(C); dot(D);
dot(E); dot(F); dot(G); dot(H);
dot(I); dot(J); dot(K); dot(L);
dot(W); dot(X); dot(Y); dot(Z);
\end{asy}
\end{wrapfigure}

\subsection{Rumors}
The following problems have claimed solutions in various levels of detail,
but all of them have some part I'm not totally happy with.
If you're an expert in cubics, please help me fill in the gaps.

\begin{problem}[\href{https://aops.com/community/p23586650}{TSTST 2021/1}]
Let $ABCD$ be a quadrilateral inscribed in a circle with center $O$.
Points $X$ and $Y$ lie on sides $AB$ and $CD$, respectively.
Suppose the circumcircles of $ADX$ and $BCY$
meet line $XY$ again at $P$ and $Q$, respectively.
Show that $OP=OQ$.
\end{problem}

\begin{problem}
  In the figure (with the outer two circles as given),
  show the four inner points are cyclic.
\end{problem}

\begin{problem}[\href{https://aops.com/community/p6580553}{TSTST 2016/6}]
Let $ABC$ be a triangle with incenter $I$, and whose incircle is tangent
to $\overline{BC}$, $\overline{CA}$, $\overline{AB}$ at $D$, $E$, $F$, respectively.
Let $K$ be the foot of the altitude from $D$ to $\overline{EF}$.
Suppose that the circumcircle of $\triangle AIB$
meets the incircle at two distinct points $C_1$ and $C_2$,
while the circumcircle of $\triangle AIC$ meets the incircle
at two distinct points $B_1$ and $B_2$.
Prove that the radical axis of the circumcircles of
$\triangle BB_1B_2$ and $\triangle CC_1C_2$ passes through the midpoint $M$ of $\overline{DK}$.
\end{problem}

\end{document}
